\section{Discussion and limitations}
\label{sec:discussion}

\paragraph{What is established.}
At fixed hidden width, widening the latent space delays memorization: the synthetic memorized fraction falls from $30\%$ to near zero over $\dlat=5\to40$, the decline survives the signal-projected test (\cref{tab:e2-projection}), and the same ordering appears as a two- to four-fold slowdown on CelebA and CIFAR-10 once the VAE preserves identity. In the frozen-feature model the delay is a ratio of spectral edges that grows because $q$ falls, and the norm-fixed control (\cref{rem:normfixed-buf}) shows that the mode count is not the driver.

\paragraph{The causal gap.}
Whether this mechanism operates in trained networks is the main open question. In the synthetic MLPs the first-layer pre-activation variance is the same at every width at initialization and at 50k steps ($q(40)/q(5)=0.96$ and $1.00$); it separates only later, growing together with memorization at small $\dlat$ (ratio $0.30$ at 5M steps). There a large $q$ is a signature of memorizing, not a width-set cause that precedes it, so the mechanism is verified in the RFNN but not shown to be what delays memorization in the MLPs. On real data the $q$-mechanism accounts for $1.7$--$1.9\times$ of the $2.1$--$3.9\times$ slowdown, roughly half in log units, and a $q$-only reading fails across datasets.

\paragraph{Boundaries.}
The spatial-DiT pilot on CelebA, with a $4\times4\times\dlat$ latent grid, does not show the decline: memorization stays high across channel counts (\cref{fig:app-spatial-dit-summary}), so the effect is established only for vector latents with MLP denoisers. All sweeps use $\nsamp\le1000$ training points, and a sweep that scales width with $\dlat$ is a capacity sweep. What would close the gap is a trained-network quantity that is set by width before memorization begins and predicts its onset; a candidate consistent with the data is the per-coordinate gap between the empirical and population scores, which grows $3.6\times$ from $\dlat=5$ to $40$.
